% =====================================================================
%  CONCLUSIONES Y RECOMENDACIONES (versión sencilla)
%  Las cifras del modelo se citan con las macros de valores_modelo.tex.
% =====================================================================
\chapter*{Conclusiones y recomendaciones}
\addcontentsline{toc}{chapter}{CONCLUSIONES Y RECOMENDACIONES}
\setcounter{section}{0}
\renewcommand{\thesection}{\arabic{section}}

\section{Conclusión general}

La investigación buscó responder de qué manera un Sistema de Gestión de la Calidad (SGC) basado en la norma ISO~9001, centrado en el control de materiales, permite reducir los costos de ejecución de las losas alivianadas con viguetas pretensadas del Edificio de Posgrado de la UATF. Como la segunda losa no se ejecutó durante el estudio, se reconstruyó con documentos la línea base de la primera losa y se evaluó el sistema \emph{ex ante}, con un modelo de costos y una simulación de Monte Carlo, sobre los 2.234,60~\mdos{} restantes.

El SGC diseñado permite reducir el costo por dos vías: disminuye el desperdicio de materiales y los retrabajos, que forman el costo de no calidad, y controla las cantidades que se miden y se pagan. Por la primera vía, la simulación estimó una reducción media del costo unitario del \rMedia{}, de Bs~\CUsinMedia{} a Bs~\CUconMedia/\mdos{}; una probabilidad de \PRpos{} de que el costo disminuya; una reducción del costo de no calidad del \CNCreduccionPct{}, y un ahorro esperado de Bs~\AhorroMedia{} en el saldo de la obra.

Sin embargo, la reducción no alcanzó la magnitud planteada en la hipótesis general. La probabilidad de que el costo unitario bajara al menos un \MCumbral{} fue de \PHG{}, muy inferior al criterio de \MCconf{}; por tanto, la hipótesis general no se acepta y se mantiene la hipótesis nula. La causa es estructural: el SGC solo actúa sobre el costo de no calidad, que en la línea base representó el \CNCsinDetPct{} del costo unitario, y el propio sistema cuesta el \CCconDetPct{}. Aun con una eficacia del 100~\%, la reducción llegaría al \rMaxTeorico{}, y alcanzar el \MCumbral{} exigiría una eficacia del \KequilibrioUmbral{}, superior a la que reporta la literatura revisada.

Las hipótesis específicas, en cambio, se aceptan: el SGC reduce el índice de desperdicio de materiales (HE1) y el costo de no calidad (HE2), con probabilidades de \PHEuno{} y \PHEdos{}, siempre que el sistema se aplique efectivamente en obra. En síntesis, el SGC reduce el costo en el sentido esperado, pero en un solo ítem su efecto es menor que el previsto. Para la UATF, una parte importante del beneficio está en asegurar la conformidad del hormigón y en controlar las cantidades que se pagan.

\section{Conclusiones específicas}

\textbf{OE1. Gestión de la calidad de la primera losa.} La primera losa se ejecutó sin un sistema de gestión de la calidad. El cumplimiento de los requisitos de la norma ISO~9001:2026 fue del 19,8~\%, con todas las cláusulas en el rango deficiente; las más débiles fueron la mejora (10,0~\%) y la planificación (12,5~\%). Se identificaron 22 no conformidades, el 54,5~\% de ellas en la medición y la gestión. Las más relevantes fueron que la losa se certificó para pago sin ensayos de resistencia de su hormigón, que los certificados de las viguetas y del EPS son posteriores al inicio del armado y que no hubo registros de consumo de materiales ni de no conformidades. Además, las cantidades de acero presentan inconsistencias documentales que deben verificar la Supervisión y la Fiscalización. La causa raíz común fue que la calidad se confió a la verificación puntual de la Supervisión antes de cada vaciado, sin criterios escritos, registros ni seguimiento (Sección~\ref{sec:diagnostico}).

\textbf{OE2. Línea base de desperdicio y de costo de no calidad.} De los cuatro materiales de la losa, solo dos tuvieron un dato documental para calcular su índice de desperdicio. Del EPS se entregaron 2.350 piezas frente a 2.237,9 teóricas, lo que da un índice de \IDMsinEps{}. Del acero se pagó un \IDMsinAce{} más que la cantidad teórica; ese valor se usó como una aproximación, porque es una cantidad pagada y no un consumo medido. Para el hormigón y las viguetas no hubo registros de consumo: para el primero se adoptó un valor de la literatura y, para las segundas, un valor supuesto, porque la literatura revisada no ofrece mediciones de referencia. Con los valores más probables, el índice de desperdicio ponderado fue de \IDMsinDet{} y el costo de no calidad, de Bs~\CNCsinDet/\mdos{}, el \CNCsinDetPct{} del costo unitario; con la incertidumbre de los parámetros, la media de ese costo fue de Bs~\CNCsinMedia/\mdos{}. El acero, los retrabajos y el hormigón concentraron la mayor parte del costo de no calidad (Sección~\ref{sec:linea-base-cuantitativa}). La falta de registros de consumo es en sí misma un resultado: sin ellos, el desperdicio de la losa no podía conocerse ni controlarse.

\textbf{OE3. Componentes del SGC.} Se estableció un SGC del proyecto organizado sobre las cláusulas 4 a 10 de la norma ISO~9001:2026 y concretado en un plan de la calidad del ítem según la norma ISO~10005:2018 (Sección~\ref{sec:diseno-sgc}). Sus componentes son la política y doce objetivos de la calidad con metas; las matrices de riesgos y de responsabilidades; un plan de puntos de inspección y ensayo (PPI) con catorce controles, cuatro de ellos puntos de espera; cinco procedimientos de control de materiales, desde la recepción hasta la conciliación de las cantidades teóricas, entregadas, colocadas, medidas y pagadas; el control del hormigón según el capítulo~26 del ACI~318-25, con al menos cinco ensayos de resistencia por losa; el procedimiento de no conformidades y acción correctiva; el registro de los costos de prevención, evaluación y fallas, y los formatos C.1 a C.8 (Anexo~\ref{anx:instrumentos}). Cada una de las nueve brechas del diagnóstico quedó vinculada con un requisito de la norma, un componente del sistema y un indicador. Con los valores más probables, el sistema cuesta Bs~\CCconDet{} por \mdos{} del saldo.

\textbf{OE4. Efecto del SGC sobre el costo y el desperdicio.} Con los valores más probables, el costo unitario bajó de Bs~\CUsinDet{} a Bs~\CUconDet/\mdos{}, una reducción del \rDet{}. En los \MCiter{} escenarios de la simulación (Sección~\ref{sec:modelo-simulacion}), la reducción media fue del \rMedia{}, con percentiles 5 y 95 de \rPcinco{} y \rPnoventaycinco{}; el índice de desperdicio ponderado medio bajó de \IDMsinMedia{} a \IDMconMedia{}, y el costo de no calidad medio, de Bs~\CNCsinMedia{} a Bs~\CNCconMedia/\mdos{}. Entre esos mismos percentiles, el ahorro en el saldo fue de Bs~\AhorroPcinco{} a Bs~\AhorroPnoventaycinco{}. Los parámetros que más influyeron en el resultado fueron el costo del SGC, la fracción de retrabajo de la línea base y la eficacia del sistema sobre los retrabajos. Con una misma eficacia para todos los materiales y los retrabajos, el SGC compensa su propio costo desde una eficacia del \KequilibrioCero{}. Si no se aplicara en obra en el \PiFalla{} de los casos, la reducción media bajaría al \rMediaFalla{}; si se incluyeran los costos ocultos de no calidad, subiría al \rMediaOcultos{}. El sistema deja de ser rentable si su eficacia es claramente menor que la reportada en la literatura o si su costo se duplica (Tabla~\ref{tab:variantes}).

\textbf{OE5. Contraste de las hipótesis y validación del SGC.} El contraste se realizó con reglas de decisión probabilísticas, equivalentes a un nivel de significancia de 0,05: una hipótesis se aceptaba si su condición se cumplía con una probabilidad de al menos \MCconf{} en los escenarios simulados (Sección~\ref{sec:contraste}). Con esa regla, la hipótesis general no se aceptó, con una probabilidad de \PHG{}, y se aceptaron HE1 y HE2, con probabilidades de \PHEuno{} y \PHEdos{}. La decisión sobre la hipótesis general se mantuvo en todas las variantes de los supuestos y en el escenario con costos ocultos, en el que la probabilidad de alcanzar el umbral fue de \PHGocultos{}. Para la validación del SGC, el estudio diseñó el instrumento de juicio de expertos: el formato C.8, en el que al menos tres ingenieros civiles con experiencia en supervisión, control del hormigón o gestión de la calidad califican la pertinencia, la claridad, la suficiencia y la aplicabilidad de cada componente del sistema, y el coeficiente V de Aiken, que acepta los componentes con un valor de al menos 0,70 (Anexo~\ref{anx:validacion}). Se recomienda aplicarlo antes de implementar el sistema en obra.

\section{Recomendaciones}

\subsection{Al Departamento de Infraestructura de la UATF}
\begin{enumerate}[label=(\arabic*)]
  \item Validar el SGC por juicio de expertos con el instrumento diseñado en este estudio (Anexo~\ref{anx:validacion}) e incorporar sus observaciones antes de exigir el sistema en la contratación del saldo de la obra.
  \item Incluir en el documento base de contratación del saldo, como obligación de la nueva contratista, un plan de la calidad del ítem con el PPI, el plan de muestreo del hormigón y los formatos del Anexo~\ref{anx:instrumentos}, y prever que el paño cuyo hormigón aún no tenga resultados a 28 días se certifique con observación y retención de pago.
  \item Antes de liquidar el contrato resuelto, disponer la conciliación elemento por elemento del cómputo del acero del ítem~8. La línea de la losa computada como 3~$\times$~4.171,18~kg y las dos versiones de la planilla de cierre son inconsistencias documentales que deben verificar la Supervisión y la Fiscalización. Además, conviene definir con el análisis de precios unitarios contractual si el ítem~13 incluye el acero de la parrilla y de los negativos.
  \item Actualizar las especificaciones técnicas del ítem: sustituir la CBH-87 y las normas antiguas de cemento, áridos y agua por la NB~1225001-1:2017 y las ediciones vigentes de IBNORCA, incluir la NB~997:2016 para las viguetas, declarar el criterio de aceptación del hormigón, definir la geometría del complemento de EPS y completar los planos con la ubicación de los nervios y de la armadura de negativos.
  \item Evaluar, con la Supervisión y el proyectista, la conveniencia de verificar en el lugar la resistencia del hormigón de la primera losa, que no tiene ensayos, antes de cargarla con los niveles superiores.
\end{enumerate}

\subsection{A las futuras empresas contratistas}
Se recomienda aplicar un SGC ligero, con el personal existente y centrado en los controles de mayor efecto: el control volumétrico del hormigón, con la producción limitada al volumen de cada vaciado; el despiece y el prearmado del acero; la recepción de los materiales por lotes, con criterios conocidos, y la conciliación de cantidades por paño. Como el costo del sistema fue uno de los parámetros de mayor influencia y su duplicación lo vuelve no rentable, conviene mantenerlo cerca del valor previsto y registrar los costos de la calidad desde el inicio, para comprobar que la reducción de las fallas supere el costo de la prevención y la evaluación. En un contrato como el estudiado, cuyos precios unitarios incluyen las reparaciones por trabajos defectuosos y no admiten reajuste, la contratista es la principal beneficiaria económica de reducir el desperdicio.

\subsection{A la Supervisión y la Fiscalización}
Se recomienda liberar cada paño en los puntos de espera del PPI, con la ficha C.2 y el asiento del Libro de Órdenes firmado por el Director de Obra; exigir la toma de probetas en cada vaciado por personal independiente de la producción y no aceptar resistencias proyectadas desde otras edades; exigir los certificados de los materiales por lote o por colada antes de su colocación; aplicar la conciliación de cantidades del procedimiento PCM-04, con la planilla de armado del acero de cada paño, antes de aprobar cada planilla de avance, y registrar las no conformidades y su cierre, para que la calificación de la calidad en los informes se apoye en evidencia objetiva.

\subsection{Para futuras investigaciones}
\begin{enumerate}[label=(\arabic*)]
  \item Medir el efecto del SGC en obra cuando se ejecute el saldo, con el protocolo de la Sección~\ref{sec:verificacion-futura}: registrar los consumos y los costos por paño homólogo, compararlos con la línea base mediante la prueba \emph{t} para muestras relacionadas o la de Wilcoxon, y actualizar con esos datos los parámetros del modelo.
  \item Incluir los costos de plazo, como las multas, los gastos generales de las ampliaciones y las esperas por material, y los costos ocultos de no calidad, que el modelo principal no valoró.
  \item Extender la evaluación a otros ítems de hormigón armado y al proyecto completo, para establecer si a esa escala un SGC alcanza reducciones del orden del \MCumbral{}.
  \item Medir el desperdicio de las viguetas pretensadas y de los complementos de EPS, para los que no se encontraron valores de referencia en la literatura revisada.
  \item Modelar las correlaciones entre los parámetros y replicar el diagnóstico y el modelo en otras obras de la UATF y de otras entidades públicas, para ampliar la validez externa.
\end{enumerate}

\section{Limitaciones del estudio}

Las conclusiones anteriores están sujetas a las limitaciones analizadas en la discusión de resultados (Sección~\ref{sec:discusion}). La evaluación del SGC fue \emph{ex ante}, porque la segunda losa no se ejecutó durante la investigación; por ello, el efecto del sistema se estimó con un modelo y no se midió en obra. Los datos provienen de una sola obra y de una sola losa, lo que limita la generalización de los resultados. La línea base no contó con registros de consumo: solo el EPS y el acero tuvieron un dato documental, y ninguno de los dos es un consumo medido; los demás parámetros se tomaron de la literatura o se supusieron. El análisis de precios unitarios contractual no estuvo disponible, por lo que se reconstruyó uno de referencia, calibrado al precio del contrato con un factor de \KappaTrece{}. Los parámetros se simularon de forma independiente, sin correlaciones entre ellos, y los costos se evaluaron a precios constantes y sin los costos de plazo. El diagnóstico frente a la norma fue documental y no incluyó entrevistas. Por último, la evaluación del SGC se basó en el modelo de costos: el instrumento de validación por juicio de expertos se diseñó, pero no se aplicó durante la investigación.

Pese a estas limitaciones, la decisión sobre la hipótesis general es robusta, porque la reducción solo supera el umbral con una eficacia cercana al 100~\%, muy superior a la que reportan los antecedentes. Las magnitudes de la reducción y del ahorro, en cambio, deben verificarse con datos medidos en obra.

% Restablece la numeración normal de secciones
\renewcommand{\thesection}{\thechapter.\arabic{section}}
